\documentclass{article}
\usepackage{../math_template}

\author{Jonathan Kasongo}
\title{Irish Math Olympiad 1988 --- P12}

\begin{document}
\maketitle

\begin{problem}{Irish Math Olympiad 1988 --- P12}
Prove that if $n$ is a positive integer, then

$$
\sum_{k=1}^n \cos^4 \left( \frac{k\pi}{2n+1} \right) = \frac{6n-5}{16}
$$
\end{problem}

\begin{solution}{Solution}
Where $z = e^{i \theta}$, we have $16\cos^4 \theta = (z + z^{-1})^4 =
(z^4 + z^{-4}) + 4(z^2 + z^{-2}) + 6$. So the sum becomes

$$
\frac{6n}{16} + \frac{1}{16}\sum_{k=1}^n \left( 2\cos(2 \phi_k) + 8\cos(\phi_k) \right)
$$

where $\phi_k = 2k\pi / (2n+1)$.\\

We notice that $e^{i \phi_0}, e^{i \phi_1}, \ldots, e^{i\phi_{2n}}$ form the
$(2n+1)$-th roots of unity. Since they solve the equation $x^{2n+1} - 1 = 0$
we deduce that they sum to 0 using Vieta's formulae. Also since
$x^{2n+1} = e^{2m \pi i k} = e^{2\pi i k}$, we get that $x = e^{mi \phi_k}
= e^{i \phi_k}$ for any positive integer $m$. Now,

\[
\begin{aligned}
\sum_{k=0}^{2n} \cos(m\phi_k) &= 0 \\
\sum_{k=1}^{2n} \cos(m\phi_k) &= -1 \\
\end{aligned}
\]

Note that $\cos(m\phi_{2n+1-k}) = \cos\left(2m\pi - \frac{2km\pi}{2n+1
}\right) = \cos(2km\pi/(2n+1)) = \cos(m \phi_k)$ and so we can write

$$
\sum_{k=1}^{2n} \cos(m\phi_k) =
\sum_{k=1}^{n} \left[\cos(m\phi_k) + \cos(m\phi_{2n+1-k}) \right]=
\sum_{k=1}^n 2\cos(m \phi_k) = -1
$$

And so our original sum becomes
$$
\frac{6n}{16} + \frac{1}{16} (-1 - 4) = \frac{6n - 5}{16}
$$

as desired. $\blacksquare$
\end{solution}

\end{document}
